
\documentclass[11pt]{article}
\usepackage[margin=1in]{geometry}

% Core packages
\usepackage{amsmath,amssymb}
\usepackage{tikz-cd}
\usepackage{multicol}
\usepackage{multirow}

% Paragraphs
\setlength{\parindent}{0pt}
\setlength{\parskip}{1\baselineskip}

\title{Práctica 6: Medidas de radiactividad ambiental}
\author{Hao Chen Miguel Ángel Montero Gómez Andrés Vinuesa Espinosa}
\date{\today}

\begin{document}
\maketitle

\section{Introducción}
\section{Metodología}
\section{Resultados}
A partir de los valores medios de cuentas en 5 minutos, se obtiene las cuentas por minuto (cpm), suponiendo
que la desintegraión radiactiva sigue una distribución de Poisson, lo cual es una hipótesis razonable, 
ya que este es un suceso poco frecuente que ocurre de forma aleatoria e indepndiente.
Se ha tomado un intervalo de confianza del 95\%,
\begin{itemize}
    \item \textbf{Cuentas por minuto (cpm)}:
    \begin{itemize}
        \item Interior: $c = \frac{103.67}{5} = \textbf{20.73} \pm \mathbf{3.99} \text{ cpm}$
        \item Exterior: $c = \frac{94.67}{5} = \textbf{18.93} \pm \mathbf{3.81} \text{ cpm}$
    \end{itemize}

    \item \textbf{Dosis de exposición en mR/h} (dividiendo las cpm entre 1000):
    \begin{itemize}
        \item Interior: $20.73 / 1000 = (0.0207 \pm 0.0040) \text{ mR/h}$
        \item Exterior: $18.93 / 1000 = (0.0189 \pm 0.0038) \text{ mR/h}$
    \end{itemize}
    
    \item \textbf{Dosis absorbida en mGy/h} ($\approx 0.01 \times \text{dosis de exposición}$):
    \begin{itemize}
        \item Interior: $0.01 \times 0.0207 = (0.000207 \pm 0.000040) \text{ mGy/h}$
        \item Exterior: $0.01 \times 0.0189 = (0.000189 \pm 0.000038) \text{ mGy/h}$
    \end{itemize}
    
    \item \textbf{Dosis absorbida diaria en mGy/día}:
    \begin{itemize}
        \item Interior: $0.000207 \times 24 = (0.00498 \pm 0.00096) \text{ mGy/día}$
        \item Exterior: $0.000189 \times 24 = (0.00454 \pm 0.00092) \text{ mGy/día}$
    \end{itemize}
    
    \item \textbf{Dosis equivalente en mSv/h} ($Q=1$ para radiación $\gamma$):
    \begin{itemize}
        \item Interior: $1 \times 0.000207 = (0.000207 \pm 0.000040) \text{ mSv/h} \quad (\approx 0.21 \pm 0.04 \, \mu\text{Sv/h})$
        \item Exterior: $1 \times 0.000189 = (0.000189 \pm 0.000038) \text{ mSv/h} \quad (\approx 0.19 \pm 0.04 \, \mu\text{Sv/h})$
    \end{itemize}
\end{itemize}
Cabe recalacar, que nuestros datos pueden estar sesgados, ya que el detector de radiación fue colocado en tres lugares puntuales arbitrarios, 
que pueden no ser una muestra representativa de la radiación de fondo en el campus, además, las medidas fueron tomadas en un día específico, lo que también puede influir a los resultados.
Una forma de mitigar este error, podría ser realizar mediciones en más lugares y en intervalos temporales más ampios, idealmente un año, o un mes de cada estación, para así 
obtener medidas que puedan tener menos sesgo experimental. Aunque a priori, la radiación de fondo debería ser relativamente constante,
si nos limitamos a asumir esta hipótesis sin realizar medidas más exhaustivas, podríamos estar cometiendo un error de muestreo, lo que podría llevar a conclusiones erróneas sobre la radiación de fondo en el campus.


\section{Comparación con fuentes de radiación conocidas}
En la siguiente tabla se muestras los valores de dosis de radiación de fuentes conocidas, tanto naturales como artificiales, que vienen dadas por CITAR EL GUIÓN,
en esta comparativa estamos asumiendo que la dosis de radiación gamma medida tanto en interior como exterior es representatvia de la radiación de fondo, y que 
en el campus no se encuentra ningún objeto radiactivo que la perturbe.

Como podemos observar, esto demuestra que la dosis de radiación gamma a la que nos exponemos por diversos procedimientos, es mucho menor
que la que recibimos de forma natural en el ambiente.

\begin{table}[h]
\centering
\renewcommand{\arraystretch}{1.2}
\setlength{\tabcolsep}{4pt}
\begin{tabular}{|l|l|c|c|c|}
\hline
\multirow{2}{*}{\textbf{Origen}} & \multirow{2}{*}{\textbf{Fuente de Radiación}} & \textbf{Dosis típica} & \multicolumn{2}{c|}{\textbf{\% de la medida respecto a la fuente}} \\
\cline{4-5}
& & \textbf{(mSv/año)} & \textbf{Interior (1.82 mSv/a)} & \textbf{Exterior (1.66 mSv/a)} \\
\hline
Artificial & Fuentes médicas (Rayos X, medicina nuclear) & 1.28 & 142.2 \% & 129.7 \% \\
Artificial & Productos de consumo (viajes, detectores, etc.) & 0.01 & 18200 \% & 16600 \% \\
Artificial & Vertidos (industria nuclear, hospitales) & $< 0.001$ & $> 182000$ \% & $> 166000$ \% \\
Natural & Radón (España media) & 1.15 & 158.3 \% & 144.3 \% \\
Natural & Alimentos y bebidas (K-40, marisco) & 0.50 & 364.0 \% & 332.0 \% \\
Natural & Rayos gamma (suelo, materiales construcción) & 0.48 & 379.2 \% & 345.8 \% \\
Natural & Radiación cósmica & 0.39 & 466.7 \% & 425.6 \% \\
\hline
\hline
\textbf{Experimental} & \textbf{Medida de radiación $\gamma$ (Interior)} & \textbf{1.82} & \textbf{100 \%} & \textbf{-} \\
\textbf{Experimental} & \textbf{Medida de radiación $\gamma$ (Exterior)} & \textbf{1.66} & \textbf{-} & \textbf{100 \%} \\
\hline
\end{tabular}
\caption{\footnotesize Comparativa de la dosis de radiación medida experimentalmente frente a fuentes comunes. El porcentaje muestra la proporción del valor experimental respecto a dichas fuentes.}
\label{tab:comparativa_gamma}
\end{table}

\end{document}
